\section{Ba kỹ thuật dùng chung}
\label{sec:ch07-ba-ky-thuat}

\index{phép nhiễu}%
\index{xấp xỉ tuyến tính}%
Mục trước đã tách ba nhánh theo đối tượng mà chúng gán điểm. Zhang và cộng sự
chồng lên đó một trục thứ hai, là kỹ thuật dùng để tính điểm, và trục này chỉ có
ba giá trị cho cả ba nhánh~\autocite{p18unified}.

Kỹ thuật thứ nhất là phép nhiễu. Ta thay đổi hoặc bỏ hẳn một đối tượng, chạy lại
mô hình, rồi lấy chênh lệch đầu ra làm điểm số. Chương~\ref{ch:shap} đã gặp dạng
tinh vi của nó khi \tn{giá trị Shapley}{Shapley value} lấy trung bình đóng góp
biên trên mọi liên minh đặc trưng. Kỹ thuật thứ hai là gradient. Thay vì thử
từng thay đổi, ta lấy đạo hàm của đầu ra theo đối tượng cần chấm điểm, nên chi
phí giảm xuống còn vài lượt truyền xuôi và ngược; đây là con đường của
chương~\ref{ch:attribution-theo-gradient}. Kỹ thuật thứ ba là \tn{xấp xỉ tuyến
tính}{linear approximation}. Ta khớp một mô hình tuyến tính vào hành vi của $f$
quanh điểm cần giải thích, rồi đọc hệ số của mô hình ấy làm điểm số, đúng như
LIME đã làm trong chương~\ref{ch:lime}.

Bảng phân loại của bài báo xếp các phương pháp tiêu biểu theo hai chiều: ba đối
tượng theo cột, ba kỹ thuật theo dòng, và mỗi kỹ thuật còn chia tiếp thành vài
biến thể~\autocite{p18unified}. Hình~\ref{fig:ch07-ba-doi-tuong} nén bảng ấy về
đúng ba dòng, mỗi ô giữ lại một hoặc hai phương pháp đại diện. Ở mức nén này cả
chín ô đều có tên phương pháp, nghĩa là mỗi kỹ thuật đã được thử trên cả ba
đối tượng.

\begin{figure}[htbp]
  \centering
  \begin{tikzpicture}[
  every node/.style={font=\small, align=center},
  head/.style={font=\small\bfseries, align=center},
  cell/.style={draw=black!55, rounded corners=2pt, fill=black!4,
    minimum width=33mm, minimum height=11mm, font=\footnotesize},
  rowlab/.style={font=\small, align=right, text width=26mm}
]
  \node[head] (cfa) at (0,0) {Đặc trưng\\của đầu vào};
  \node[head, right=4mm of cfa] (cda) {Dữ liệu\\huấn luyện};
  \node[head, right=4mm of cda] (cca) {Thành phần\\trong mô hình};

  \node[cell, below=6mm of cfa] (p1) {Occlusion, SHAP};
  \node[cell, right=4mm of p1] (p2) {Leave-one-out,\\Data Shapley};
  \node[cell, right=4mm of p2] (p3) {Vá kích hoạt};

  \node[cell, below=4mm of p1] (g1) {Integrated Gradients};
  \node[cell, right=4mm of g1] (g2) {Hàm ảnh hưởng,\\TracIn};
  \node[cell, right=4mm of g2] (g3) {Attribution patching};

  \node[cell, below=4mm of g1] (l1) {LIME};
  \node[cell, right=4mm of l1] (l2) {Datamodel, TRAK};
  \node[cell, right=4mm of l2] (l3) {COAR};

  \node[rowlab, left=4mm of p1.west, anchor=east] {Phép nhiễu};
  \node[rowlab, left=4mm of g1.west, anchor=east] {Gradient};
  \node[rowlab, left=4mm of l1.west, anchor=east] {Xấp xỉ\\tuyến tính};
\end{tikzpicture}

  \caption{Ba kỹ thuật theo dòng, ba đối tượng theo cột, và một phương pháp đại
  diện trong mỗi ô. Lưới là bản nén của sách, gộp tám dòng chi tiết trong bảng
  phân loại của bài báo về ba kỹ thuật~\autocite{p18unified}; các phương pháp
  trong ô lấy từ bảng đó.}
  \label{fig:ch07-ba-doi-tuong}
\end{figure}

Chỗ trống chỉ lộ ra khi trải bảng về đủ tám dòng của bài báo, tức khi mỗi kỹ
thuật được tách thành các biến thể của nó, và chỗ trống ấy cũng mang thông tin. Giá
trị Shapley đã được dùng cho cả ba đối tượng, nhưng những khái niệm khác của lý
thuyết trò chơi hợp tác thì phần lớn còn nằm trong attribution đặc trưng: mới
một khái niệm sang tới \tn{attribution dữ liệu}{data attribution}, và chưa cái
nào sang \tn{attribution thành phần}{component attribution}. Ma trận Hessian
cũng vậy, đã vào attribution đặc trưng qua Integrated Hessians và vào
attribution dữ liệu qua họ \tn{hàm ảnh hưởng}{influence function}, mà chưa được
thử cho attribution thành phần. Các tác giả đọc mỗi ô trống như một hướng nghiên
cứu: một ý tưởng đã kiểm chứng ở nhánh này và chưa ai mang sang nhánh
kia~\autocite{p18unified}.

Cách đọc ấy chỉ đứng vững nếu các ô thật sự chỉ cùng một thứ. Attribution
dữ liệu và attribution thành phần là hai nhánh Phần~II chưa chạm tới, nên chúng
được xem xét trước, sau đó mới đến khung hình thức chung.
